\subNsection{1}{Perspektiven}{Kristian Popov} \label{6_2_Perspektiven}

Die modulare Architektur erlaubt eine flexible Erweiterung der Anwendung. Künftige Perspektiven gliedern sich in folgende Schwerpunkte:

\begin{itemize}
    \item \textbf{Weitere Lernbereiche und Topics:} Neue Themen – etwa von anderen Professoren beigesteuerte Fachmodule – lassen sich über den zentralen Modulkatalog dynamisch in Navigation und Routing integrieren.
    \item \textbf{Aufgabenverwaltung und -varianten:} Lehrende können Aufgaben und Vorlagen direkt über die Oberfläche anlegen, automatisch Varianten erzeugen sowie Quizfragen und Lernkarten ergänzen.
    \item \textbf{Mobile Nutzung:} Optimierung der Benutzeroberfläche – insbesondere von Übergangstabellen und Zustandsdiagrammen – für kleine Bildschirme und Tablets.
    \item \textbf{Java-Zeilensperrung:} Abgrenzung bearbeitbarer Code-Bereiche im Editor durch Marker, um nicht zu bearbeitende Zeilen schreibgeschützt zu halten.
    \item \textbf{Gamification und LMS-Anbindung:} Motivationselemente wie Streaks oder Abzeichen sowie eine Moodle-Anbindung zur zentralen Lernstandsauswertung.
    \item \textbf{Offene Punkte aus dem Lastenheft (Abschnitt~\ref{5_2_LH}):}
    \begin{itemize}
        \item \textit{Usability-Erhebung:} Durchführung strukturierter Tests zur Bedienbarkeit über einen internen Testzugang.
        \item \textit{Barrierefreiheit \& Mobil-Profiling:} Automatisierte CI-Prüfungen für Kontraste, Tastaturbedienung und mobile Performance-Messungen.
    \end{itemize}
\end{itemize}

Während sich neue Lernbereiche, Aufgabenvarianten und Editor-Sperren vollständig in die bestehende, wartungsfreie Client-Architektur integrieren lassen, würden zentrale Auswertungen, Nutzerkonten oder eine direkte Moodle-Anbindung ein Backend erfordern.